\section{Conclusion}
\label{sec:conclusion}

This paper proposes frequent losers---firms that never win yet
participate in abnormally many public tenders---as a screening marker
for cover bidding in procurement auctions. Using 4.5 million
tender-items from Brazil's BEC platform (2009--2019), we find strong
support for FL as a screening marker (H1): FL-present tenders exhibit
3.6--8.0\% higher prices across five non-IV estimators, FL firms
co-participate with CADE-convicted
cartelists at 3.5 times the expected rate ($p < 0.001$), and the FL
screen achieves AUC~$= 0.94$---dominating bid-level screens
(AUC~$= 0.79$) while requiring only participation data.

The evidence also supports cover bidding as the causal mechanism (H2):
the FL--price effect concentrates in competitive markets
(Proposition~\ref{prop:market_selection}), Bajari--Ye tests reject
exchangeability and conditional independence, event-study evidence
reveals strategic market selection preceding cover-bidder deployment,
and structural estimation selects Regime~2
($\hat{\sigma}_c / \hat{\sigma}_g = 0.72$), indicating coordinated
but imperfectly controlled cover-bid placement. Removing the convite
minimum-bidder rule would save R\$74M--R\$211M, and the
welfare-maximizing screening threshold (1.6$\times$IQR) is close to the
baseline (1.5$\times$).

\paragraph{Caveats.}
The staggered DiD is underpowered; the FL definition necessarily
involves misclassification; and strategic adaptation (rotating cover
bidders, allowing occasional wins) is a concern for operational
deployment. The network-split results present an unresolved puzzle:
concentrated-market FL firms---those with the strongest network
signatures of coordination---show no significant price effect. One
interpretation is that cartels in concentrated markets already possess
sufficient market power to sustain supra-competitive prices without
cover bidding; FL firms there serve as redundant insurance rather than
active price inflators. Distinguishing this from alternative
explanations (e.g., compositional differences across product markets)
requires richer data on cartel structure.

\paragraph{Magnitude in context.}
The 3.6--8.0\% markup range falls at the low end of the cartel
overcharge distribution documented by \citet{connor2007price}, who
report a median of 23\% across 674 cartels. This is consistent with
our identifying a \emph{screening signal} rather than a full cartel
overcharge: FL presence marks tenders where cartels operate, but
the fixed-effects specification absorbs within-item price variation
unrelated to FL, and the cross-fit estimate (3.6\%) further
discounts any in-sample mechanical correlation. The true cartel
markup likely exceeds the FL coefficient, which captures only the
incremental price effect associated with cover-bidder deployment.

\paragraph{Policy implications.}
The FL screen requires only participation and outcome data, produces
firm-level flags for targeted investigation, and complements bid-level
screens in a two-stage workflow. Estimated welfare loss ranges from
R\$400M (cross-fit) to R\$2.4B (IV) over 2009--2019. FL flags are a
starting point for investigation, not standalone enforcement evidence.

The screen complements rather than substitutes for leniency programs.
Leniency targets insiders willing to defect; FL screening targets
the observable footprint of cover-bidding operations from the outside.
In jurisdictions where leniency uptake is low---as in much of Latin
America---proactive screening may be the only feasible detection
channel. Conversely, FL flags can inform leniency outreach by
identifying firms most likely involved in active cartels.

\paragraph{Portability.}
The FL screen's data requirements---participation records and
outcomes---are available in most electronic procurement systems.
The approach is portable to any setting where cover bidding
generates repeat-losing firms: other Brazilian states, EU member
states with centralized e-procurement (TED), and developing economies
adopting digital procurement platforms. The key institutional
requirement is that participation records include firm identifiers
and outcome flags; bid-level data are not needed.

\paragraph{Future research.}
Four directions merit investigation: external validation beyond BEC
(other states, other countries, federal procurement);
deeper network analysis (shared representatives, financial links);
dynamic models of FL-screen--leniency interaction; and field
experiments in which procurement agencies adopt FL screening.
